\chapter{Serverless ETL with AWS Glue}
\index{AWS Glue}\index{DynamicFrame}\index{Glue Studio}\index{Glue DataBrew}\index{job bookmarks}\index{DPU}\index{serverless ETL}%
\begin{objectives}
\begin{itemize}
    \item Explain Glue's serverless ETL model and the difference between DynamicFrames and Spark DataFrames.
    \item Build visual ETL pipelines in Glue Studio and profile data with Glue DataBrew.
    \item Make re-runs safe with job bookmarks and idempotent write patterns.
    \item Tune jobs with worker types, DPU sizing, and Flex execution, and reason about DPU billing.
    \item Evaluate Glue versus EMR for analyst-built daily reporting pipelines.
    \item Build a pipeline that joins two tables, drops nulls, transforms timestamps, and loads clean data into Redshift.
\end{itemize}
\end{objectives}

\begin{crosscloud}
Serverless, catalog-integrated ETL is a universal pattern; only the branding differs.

\vspace{2mm}
{\footnotesize
\begin{tabular}{@{}p{2.8cm}p{3.0cm}p{3.0cm}p{3.0cm}@{}}
\toprule
\textbf{Capability} & \textbf{AWS} & \textbf{Azure} & \textbf{GCP} \\
\midrule
Serverless ETL engine & AWS Glue (Spark) & Fabric Dataflow Gen2 / pipelines & Dataflow / Cloud Data Fusion \\
Visual ETL designer & Glue Studio & Data Factory designer & Data Fusion studio \\
Interactive data prep & Glue DataBrew & Power Query (Fabric) & Dataprep (Trifacta) \\
Shared catalog & Glue Data Catalog & OneLake catalog / Purview & Dataplex \\
Re-run state & Job bookmarks & Pipeline state / watermark columns & Job state / watermarking \\
\addlinespace
Capacity unit & DPU-hour & Capacity units & vCPU-hr (Dataflow) \\
Batch spot discount & Flex execution & Low-priority pools & Preemptible/Spot workers \\
\bottomrule
\end{tabular}}

\vspace{1mm}
\textbf{Snowflake} approaches the same problems from inside the warehouse: streams, tasks, and Snowpark replace external ETL for in-warehouse transforms. \textbf{MongoDB} users meet this layer as the Atlas Data API or Spark connector feeding the lake. The Glue skill set---schema-aware transforms, idempotent re-runs, capacity tuning---transfers to all of them.
\end{crosscloud}

\section{Week 10 Roadmap and Mental Model}
Chapter~9 gave you full control of a big-data cluster; Week~10 asks when you should decline that control. AWS Glue is serverless Spark (and Python shell) with the Data Catalog built in: you write the transformation, Glue provisions the workers, runs the job, and tears everything down \cite{awsGlueDocs}. For the large middle class of ETL work---daily joins, type fixes, format conversion, Redshift loads---Glue removes cluster operations entirely.

The week's three engineering disciplines: understand the \emph{abstraction} (DynamicFrames and when to drop to DataFrames), make every job \emph{idempotent} (bookmarks plus overwrite-style writes), and manage \emph{capacity economics} (DPUs, worker types, Flex).

\begin{definitionbox}
\textbf{AWS Glue} is a serverless data integration service comprising an ETL engine (Spark and Python shell), a visual designer (Glue Studio), an interactive preparation tool (DataBrew), a scheduler and workflow engine, crawlers, and the Data Catalog shared with Athena, Spectrum, and EMR.
\end{definitionbox}

\begin{keyterms}
\textbf{DynamicFrame}: Glue's schema-flexible distributed collection with resolve-choice semantics. \textbf{DPU}: data processing unit---4 vCPU + 16 GB; the billing and capacity unit. \textbf{Job bookmark}: persisted per-job state tracking processed inputs so re-runs process only new data. \textbf{Glue Studio}: the visual job designer that generates PySpark. \textbf{DataBrew}: the interactive, no-code profiling and cleaning tool. \textbf{Flex}: discounted spare-capacity execution for non-urgent jobs.
\end{keyterms}

\section{Monday: Glue ETL---DynamicFrames Versus Spark DataFrames}
\subsection{The DynamicFrame Abstraction}
Glue jobs are Spark jobs with an extra abstraction. A Spark \emph{DataFrame} has one rigid schema; a \emph{DynamicFrame} is a collection of self-describing records designed for data whose schema is messy or evolving \cite{awsGlueDocs}. Its signature operations---\texttt{resolveChoice}, \texttt{relationalize}, \texttt{applyMapping}---handle the real-world cases: a field that is sometimes a string and sometimes a struct, nested JSON that must become relational tables.
\index{DynamicFrame}\index{resolveChoice}

\begin{lstlisting}[language=Python,caption={Glue job skeleton: catalog in, DynamicFrame transforms, Redshift out.},label={lst:ch10-glue-job}]
import sys
from awsglue.context import GlueContext
from awsglue.dynamicframe import DynamicFrame
from awsglue.job import Job
from awsglue.utils import getResolvedOptions
from pyspark.context import SparkContext
from pyspark.sql.functions import col, to_timestamp

args = getResolvedOptions(sys.argv, ["JOB_NAME"])
sc = SparkContext()
glueContext = GlueContext(sc)
job = Job(glueContext)
job.init(args["JOB_NAME"], args)

orders = glueContext.create_dynamic_frame.from_catalog(
    database="week7_lake", table_name="sales_parquet",
    transformation_ctx="orders",
)
customers = glueContext.create_dynamic_frame.from_catalog(
    database="week7_lake", table_name="customers",
    transformation_ctx="customers",
)

# Drop to DataFrames for the join/filter/timestamp logic...
joined = (
    orders.toDF()
          .join(customers.toDF(), "customer_id")
          .filter(col("customer_id").isNotNull() & col("amount").isNotNull())
          .withColumn("order_ts", to_timestamp(col("order_date")))
)

# ...and back to a DynamicFrame for Glue's Redshift sink options.
out = DynamicFrame.fromDF(joined, glueContext, "out")

glueContext.write_dynamic_frame.from_options(
    frame=out,
    connection_type="redshift",
    connection_options={
        "url": "jdbc:redshift://week5-wg.123456789012.us-east-1.redshift-serverless.amazonaws.com:5439/dev",
        "dbtable": "analytics.sales_clean",
        "redshiftTmpDir": "s3://acme-analytics-raw-123456789012/tmp/glue/",
        "aws_iam_role": "arn:aws:iam::123456789012:role/RedshiftS3ReadRole",
    },
    transformation_ctx="redshift_sink",
)
job.commit()
\end{lstlisting}

\textbf{When to use which:} start from the catalog with DynamicFrames, convert to DataFrame (\texttt{.toDF()}) the moment you need the full Spark API---window functions, complex expressions, ML---and convert back for Glue-specific sinks. Most production jobs live mostly in DataFrame land; the DynamicFrame earns its keep at schema boundaries.
\index{DataFrame conversion}

\section{Tuesday: Glue Studio and DataBrew}
\subsection{Visual ETL for Working Teams}
Glue Studio lets you compose sources, transforms, and targets on a canvas and generates the PySpark job underneath. Its honest role in a professional estate: fast prototyping, and a legible starting point for analysts. The generated code should graduate into version control, where engineers harden it---parameterize paths, add error handling, attach bookmarks.
\index{Glue Studio}

\subsection{DataBrew: Profiling Before Pipelines}
Glue DataBrew is the interactive counterpart: point it at a dataset and it profiles columns (null rates, distributions, outliers) and offers 250+ no-code transformations as a visual ``recipe'' that can run on a schedule \cite{awsGlueDataBrewDocs}. DataBrew's highest-leverage use is \emph{before} you write a Glue job---its profile tells you the null rates and type inconsistencies your job must handle. Chapter~12's capstone uses DataBrew again as a data-quality gate.
\index{Glue DataBrew}\index{data profiling}

\begin{tipbox}
DataBrew profile jobs scan the full dataset and bill accordingly. Profile a sampled S3 path during development; reserve full-dataset profiles for the pre-production validation run.
\end{tipbox}

\section{Wednesday: Bookmarks, DPUs, and Flex}
\subsection{Job Bookmarks and Idempotency}
A job bookmark records what a job already processed---which S3 objects, which partition values---so a re-run processes only what is new \cite{awsGlueDocs}. Bookmarks solve the ``what's new'' half of idempotency; the other half is the \emph{write} side. A job that appends rows to the target will duplicate them on any re-run that reprocesses input. Production discipline: write with overwrite semantics scoped to the partition being rebuilt (or MERGE into Redshift), parameterize the run date, and keep bookmarks enabled so retries reprocess only new partitions.
\index{job bookmark}\index{idempotency}

\begin{pitfall}
\textbf{Bookmark on, append write.} Bookmarks prevent re-reading; they do nothing about re-writing. A mid-job failure retried with an append sink duplicates everything the first attempt wrote. Pair bookmarks with partition-overwrite or merge sinks, and make the pair a code-review checklist item.
\end{pitfall}

\subsection{Capacity and Cost: DPUs, Worker Types, Flex}
A DPU is 4 vCPU and 16 GB of memory; jobs bill per DPU-second with a one-minute minimum \cite{awsGlueDocs}. Tuning levers: worker type (G.1X, G.2X and larger for memory-heavy shuffles), worker count (parallelism versus diminishing returns), and \emph{Flex execution}, which runs the job on spare capacity at a discount in exchange for variable start times---ideal for non-urgent backfills and nightly jobs with slack in their window.
\index{DPU}\index{worker type}\index{Flex execution}

\section{Thursday Hands-on Project: A Glue Studio Pipeline into Redshift}
\index{hands-on project!Glue Studio pipeline}
Build a visual ETL pipeline in Glue Studio that joins two tables, drops null values, transforms a timestamp, and loads the clean data into Redshift---then harden the generated script.

\subsection*{Lab Objectives}
\begin{itemize}
    \item Compose the pipeline visually: two catalog sources, a join, a null filter, a timestamp transform, a Redshift target.
    \item Inspect and parameterize the generated PySpark script.
    \item Enable bookmarks, run the job twice, and prove the second run loads nothing new.
    \item Introduce a bad record and observe the pipeline's behavior.
\end{itemize}

\subsection*{Procedure}
\begin{enumerate}
    \item In Glue Studio, create a job from the visual canvas. Add two \emph{AWS Glue Data Catalog} sources: \texttt{week7\_lake.sales\_parquet} and \texttt{week7\_lake.customers}.
    \item Add a \emph{Join} transform on \texttt{customer\_id}; a \emph{Filter} transform dropping rows where \texttt{customer\_id IS NULL} or \texttt{amount IS NULL}; a \emph{Custom Transform} (or SQL) converting \texttt{order\_date} to a true timestamp column \texttt{order\_ts}.
    \item Add an \emph{Amazon Redshift} target writing \texttt{analytics.sales\_clean} with the S3 temporary directory set to a dedicated \texttt{tmp/glue/} prefix.
    \item Open the generated script tab. Replace hardcoded bucket and table literals with job parameters; save.
    \item On the \emph{Job details} tab: worker type G.1X, 4 workers, job bookmarks \emph{Enabled}, Flex \emph{off} for the timed first run.
    \item Run the job; verify row counts in Redshift. Load one new day's files to S3, rerun, and verify only the new partition's rows arrive.
\end{enumerate}

\begin{lstlisting}[language=SQL,caption={Verification queries against the Redshift target.},label={lst:ch10-verify}]
-- Row count and freshness
SELECT COUNT(*) AS rows_loaded,
       MAX(order_ts) AS latest_order
FROM analytics.sales_clean;

-- After the second (bookmark) run, this delta must equal only new rows:
SELECT order_date::DATE, COUNT(*)
FROM analytics.sales_clean
GROUP BY 1
ORDER BY 1 DESC
LIMIT 5;

-- Null hygiene: both must return zero
SELECT COUNT(*) FROM analytics.sales_clean
WHERE customer_id IS NULL OR amount IS NULL;
\end{lstlisting}

\subsection*{Verification}
\begin{itemize}
    \item The Redshift table contains exactly the joined, filtered rows---no nulls in the enforced columns.
    \item The second run's Redshift delta equals the new data only, demonstrating bookmark behavior plus idempotent writes.
    \item The generated script is under version control with parameters externalized.
    \item A deliberately malformed record produces a documented, non-silent outcome (dropped by filter, or job failure---your stated choice).
\end{itemize}

\section{Friday Use Case: Glue or EMR for the Analyst Team?}
\index{use case!Glue versus EMR}
A team of eight business analysts must build and run daily reporting pipelines: join the day's extracts, clean types, load Redshift. None manage infrastructure; two know Python well. Their director asks whether to standardize on Glue or EMR.

\subsection{The Evaluation}
\textbf{Skill fit.} Glue Studio's visual canvas and generated code match analysts who think in transforms, not clusters. EMR's value---framework choice, instance tuning, long-running jobs---is precisely what this team does not want to operate.

\textbf{Workload fit.} The pipelines are daily, minutes-to-an-hour, schema-evolving, and catalog-centric: the Glue sweet spot. None are multi-framework, none run 24/7, none need exotic libraries beyond what custom Python wheel uploads provide.

\textbf{Cost fit.} Intermittent daily jobs bill only DPU-seconds used; there is no idle cluster. Flex execution discounts the overnight schedule further.

\textbf{Where the answer would flip.} Heavy continuous streaming, sub-second latency, framework requirements Glue lacks, or a platform team that already runs EKS---those point back to Chapter~9. State them in the recommendation so the decision is revisitable.

\begin{table}[H]
\centering
\small
\begin{tabular}{@{}p{3.4cm}p{4.4cm}p{4.4cm}@{}}
\toprule
\textbf{Criterion} & \textbf{AWS Glue} & \textbf{Amazon EMR} \\
\midrule
Operations burden & None (serverless) & Cluster lifecycle, even if transient \\
Skill required & SQL/Python, visual design & Spark/Hadoop operations \\
Workload shape & Batch, intermittent, catalog-centric & Continuous, multi-framework, exotic \\
Cost shape & DPU-seconds per run & Node-hours per cluster \\
Escape hatches & Drop to DataFrame; custom wheels & Anything the ecosystem offers \\
\bottomrule
\end{tabular}
\caption{Glue versus EMR for the analyst reporting estate.}
\label{tab:ch10-glue-vs-emr}
\end{table}

\begin{pitfall}
\textbf{Choosing the tool the builders cannot operate.} The correct comparison is not feature lists; it is who owns a 2~a.m.\ failure. A serverless pipeline an analyst can rerun beats an elegant cluster only an absent engineer can restart.
\end{pitfall}

\section{Interview Q\&A}
\subsection{Concepts and Trade-offs}
\begin{enumerate}
    \item \textbf{Q: What does a DynamicFrame offer that a Spark DataFrame does not?}\\
    A: Self-describing records with schema-resolution operators---\texttt{resolveChoice}, \texttt{applyMapping}, \texttt{relationalize}---built for messy or evolving schemas and catalog-aware sinks. Convert to DataFrame when you need the full Spark API.
    \item \textbf{Q: What do job bookmarks guarantee on a re-run?}\\
    A: That inputs already processed are not reprocessed---new objects or partitions only. They do not make the sink idempotent; overwrite or merge writes must cover the output side.
    \item \textbf{Q: When is Flex execution the cheaper choice?}\\
    A: For non-urgent jobs---backfills, overnight batch with slack---that tolerate delayed starts. Flex runs on spare capacity at a discount; latency-sensitive or SLA-bound jobs stay on standard execution.
    \item \textbf{Q: Why is Glue usually a better fit than EMR for analyst-built pipelines?}\\
    A: No clusters to operate, a visual designer that generates real code, catalog integration, and per-second billing that matches intermittent daily runs.
    \item \textbf{Q: What compute does one DPU represent for billing?}\\
    A: 4 vCPU and 16 GB of memory, billed per second with a one-minute minimum. Worker type and count multiply from there.
    \item \textbf{Q: What is DataBrew's best use in a Glue estate?}\\
    A: Interactive profiling and recipe prototyping before pipeline development---and, as Chapter~12 shows, a schedulable data-quality gate inside an orchestrated pipeline.
\end{enumerate}

\section{Further Reading}
The Glue Developer Guide covers jobs, DynamicFrames, bookmarks, worker types, and Flex \cite{awsGlueDocs}; the DataBrew guide covers profiling and recipes \cite{awsGlueDataBrewDocs}. Spark's own documentation underpins everything the DataFrame half of a Glue job does \cite{apacheSparkProject,zaharia2016apache}. The Well-Architected Analytics Lens frames the Glue-versus-EMR operational trade \cite{awsWellArchitectedAnalytics}.

\section*{Key Takeaways}
\begin{itemize}
    \item Glue is serverless Spark plus the catalog: the right default for intermittent, schema-evolving batch ETL.
    \item DynamicFrames handle messy schemas at the boundary; DataFrames carry the serious transform logic.
    \item Idempotency is a pair: bookmarks for the read side, overwrite-or-merge sinks for the write side, plus a run-date parameter.
    \item DPU sizing, worker types, and Flex execution are the cost levers; per-second billing rewards right-sizing.
    \item Glue Studio accelerates building; DataBrew accelerates understanding the data first.
    \item Tool selection is an operations decision: match the tool to the team that must run it at 2~a.m.
\end{itemize}

\section{Exercises}
\begin{enumerate}
    \item \textbf{(Conceptual)} Explain why bookmarks alone do not prevent duplicate rows in Redshift, and name the two write-side patterns that complete idempotency.
    \item \textbf{(Conceptual)} When would a Flex job be the wrong choice even if it is cheaper?
    \item \textbf{(Hands-on)} Complete the Thursday lab; then modify the script to MERGE into Redshift by business key instead of inserting, and demonstrate a safe full rerun.
    \item \textbf{(Hands-on)} Rerun the job on Flex three times and record start latency versus the standard runs.
    \item \textbf{(Hands-on)} Profile the source dataset in DataBrew; list three findings that change the job's transform logic.
    \item \textbf{(Design)} A job's runtime doubled after the source added ten columns. Walk through the DPU, worker-type, and layout checks you would run.
    \item \textbf{(Design)} Design the parameter contract (run date, source prefix, target table, mode) for a reusable nightly Glue job template.
    \item \textbf{(Architecture)} The analyst team now needs a five-minute micro-batch. Evaluate Glue streaming, EMR, and Kinesis/Flink (Chapters 13--15) for this requirement.
\end{enumerate}

\section{Capstone Project: An Idempotent Nightly ETL Fleet on Glue}\label{sec:ch10-capstone}
\index{capstone project}\index{AWS Glue!capstone}

\begin{capstonebox}
\textbf{Mission.} Convert three manual analyst workflows into a governed, idempotent nightly Glue fleet: catalog-driven sources, bookmarked incremental loads, Flex scheduling, and Redshift targets---with proof that failures rerun safely.

\textbf{Estimated time.} 4--6 hours in a sandbox AWS account.

\textbf{What you\textquotesingle{}ll practice.}
\begin{itemize}
    \item Graduating Glue Studio prototypes into version-controlled, parameterized jobs.
    \item Bookmark plus overwrite-sink idempotency, demonstrated under forced failure.
    \item Capacity tuning: worker types, DPU counts, Flex versus standard execution.
    \item Redshift sink hygiene: temporary directories, IAM roles, and merge semantics.
\end{itemize}
\end{capstonebox}

\subsection*{Scenario}
Three analyst-built nightly pipelines currently run as manual notebook executions. They overlap in sources, drift in logic, and once a month someone reloads a full table by accident. Standardize them as Glue jobs with a shared template, and prove the failure story before finance relies on the numbers.

\subsection*{Requirements}
\textbf{Functional requirements.}
\begin{itemize}
    \item Three jobs from one parameterized template: run date, source, target, and mode as parameters.
    \item Bookmarks enabled; sinks use partition-overwrite or merge, never blind append.
    \item Nightly schedule on Flex where the window allows; alerting on failure to SNS.
    \item A forced mid-job failure followed by a rerun that produces zero duplicates.
\end{itemize}

\textbf{Non-functional requirements.}
\begin{itemize}
    \item All job code in version control; no console-only configuration.
    \item DPU sizing justified with a before/after measurement for at least one job.
    \item A runbook entry per job: what it does, how to rerun safely, what the alarm means.
\end{itemize}

\subsection*{Implementation Steps}
\begin{enumerate}
    \item Extract the common template from the three workflows; externalize parameters.
    \item Implement the sinks with overwrite/merge semantics; enable bookmarks.
    \item Schedule on Flex; wire failures to SNS.
    \item Kill one job mid-run (stop the run) and rerun; capture the zero-duplicate evidence.
    \item Measure one job across two worker types and document the sizing decision.
    \item Write the runbook entries and the template README.
\end{enumerate}

\subsection*{Deliverables}
\begin{itemize}
    \item The template and three job configurations in version control.
    \item The forced-failure/rerun evidence with row-count comparisons.
    \item The sizing measurement and Flex-versus-standard latency notes.
    \item SNS alarm configuration evidence and the three runbook entries.
\end{itemize}

\subsection*{Acceptance Criteria}
\begin{itemize}
    \item Any job can be rerun for any date with no duplicates and no manual cleanup.
    \item Failure notifications arrive in SNS with job name and run ID.
    \item No pipeline step exists only in the console.
    \item Cost per nightly run is known and recorded per job.
\end{itemize}

\subsection*{Stretch Goals}
\begin{itemize}
    \item Add a Glue workflow chaining the three jobs with a shared failure path.
    \item Introduce a DataBrew profile job as a pre-load quality gate (preview of Chapter~12).
    \item Add Glue Data Quality rules to one job (deep coverage in Chapter~22).
    \item Parameterize the template enough to onboard a fourth pipeline in under an hour; demonstrate it.
\end{itemize}
